\SourceLeaf{562}Falls darauf recurriren zu. können, aufbewahrt werden müssen. Für die Zeiten, wo ungewöhnlich starke Bewegungen vorkommen, bleibt jedoch die sofortige vollständige Mittheilung wünschenswerth. Ausser den Beobachtungszahlen sind die sonstigen damit in Verbindung stehenden Umstände, der Werth der Scalentheile (oder die Messungen, auf denen die Bestimmung beruht), die Schwingungsdauer, Stand und Gang der Uhr, Namen der Beobachter, Erläuterungen zu solchen Beobachtungen, die etwa als zweifelhaft bezeichnet werden, u.\ dergl.\ beizufügen. Dass endlich immer eine \textit{baldige} Einsendung gewünscht werden muss, bedarf keiner Erinnerung.

\Needspace{24\baselineskip}\vspace{6mm}
\pdfbookmark[1]{Auszug aus dreijährigen täglichen Beobachtungen der magnetischen Declination zu Göttingen.}{dreijaehrige-declination}
\begin{center}\large\textit{Auszug aus dreijährigen täglichen Beobachtungen der magnetischen Declination zu Göttingen.}\end{center}
\begin{center}\small Resultate aus den Beobachtungen des magnetischen Vereins.\quad 1836.\quad III.\end{center}

Bei dem unaufhörlichen Wechsel kleinerer und grösserer Schwankungen in der magnetischen Declination, die wir unregelmässige nennen, insofern ihr Vorkommen an keine Zeitregel gebunden ist, gibt es zum Ausscheiden des Regelmässigen keinen andern Weg, als eine grosse Menge von Beobachtungen nach einem bestimmten Plane anzustellen, mit beharrlicher Consequenz eine lange Zeit fortzusetzen, und in schicklichen Combinationen Mittelwerthe abzuleiten, aus welchen der Einfluss der das Einzelne stets treffenden Anomalien, so viel zu erreichen möglich ist, verschwindet. Während der Vormittagsstunden nimmt in unsern Gegenden die Declination gewöhnlich zu, aber einen Tag viel, einen andern wenig, ja zuweilen (wenn auch selten) beobachtet man in der Stunde, wo gewöhnlich die Declination am grössten ist, eine kleinere, als in den Frühstunden desselben Tages. Die Ursache der vormittägigen Zunahme mag immerhin an jedem Tage wirksam sein: aber die Wirkung wird durch andere regellos dazwischen kommende Kräfte zuweilen vergrössert, zuweilen vermindert, zuweilen ganz verdunkelt. Wie viel also eigentlich die regelmässige Ursache wirkt, wie sie in den verschiedenen Jahreszeiten ungleich wirkt, lässt sich nicht aus einzelnen oder wenigen Tagen, sondern nur durch Mittelwerthe aus sehr vielen Tagen
\SourceLeaf{563}erkennen. Auf ähnliche Weise verhält es sich mit den allmählich, aber wenigstens auf sehr lange Zeit in einerlei Sinn fortschreitenden Änderungen, die wir säculare nennen, weil ihre Anhäufung auf viele Grade eine lange Reihe von Jahren erfordert. Einzelne Beobachtungen, die nur einige wenige Jahre von einander entfernt sind, mögen sie immerhin an einerlei Monatstag und zu gleicher Stunde angestellt sein, können uns darüber noch gar keine sichere Belehrung geben: aber consequent gewonnene Mittelzahlen lassen uns das schon nach wenigen Jahren anticipiren, was sonst mit einiger Annäherung erst nach mehrern Jahrzehnden festgestellt werden könnte.

Von diesem Gesichtspunkt ausgehend habe ich unter die im hiesigen magnetischen Observatorium anzustellenden Beobachtungen gleich vom Anfang an die tägliche Bestimmung der absoluten Declination, immer zu denselben Stunden, mit aufgenommen. Um jedoch leichter auf die Thunlichkeit einer langen und ununterbrochenen Fortsetzung rechnen zu können, wodurch Arbeiten dieser Art erst ihren Werth erhalten, habe ich lieber zuerst einen beschränkten Plan wählen, als auf einmal zu viel umfassen wollen. Deshalb werden täglich nur zwei Bestimmungen gemacht, Vormittags um 8 Uhr, und Nachmittags um 1 Uhr nach mittlerer Zeit. Diese mit andern Obliegenheiten am leichtesten vereinbare Stundenwahl empfahl sich auch dadurch, dass bei einem regelmässigen Verlauf der magnetischen Bewegungen der Stand der Nadel um 1 Uhr Nachmittags immer wenig von dem Maximum der Declination, so wie um 8 Uhr Vormittags in dem grössern Theile des Jahres wenig von dem Minimum entfernt ist. Das Beobachten zu bestimmten Stunden \textit{wahrer} Sonnenzeit wäre allerdings an sich noch etwas mehr naturgemäss gewesen, allein die Rücksicht auf die viel grössere Bequemlichkeit einer Anordnung nach mittlerer Zeit musste hier, wo es hauptsächlich nur auf eine consequente Durchführung nach einerlei Princip ankam, überwiegen.

Diese regelmässigen Aufzeichnungen haben mit dem ersten Januar 1834 den Anfang genommen: indessen sind die ersten drittehalb Monate von dem folgenden Auszuge ausgeschlossen, weil während dieser Zeit öfters nöthig gewordene Aufwindungen des Aufhängungsfadens Veränderungen des Nullpunkts der Torsion hervorgebracht hatten, die anfangs nicht genug beachtet wurden. Vom 17. März an ist ein stärkerer (zweihundertfacher) Aufhängungsfaden gebraucht, nachdem dessen Torsions-Nullpunkt vorher genau berichtigt war; so oft später eine Veränderung mit diesem Faden oder in Beziehung auf einen andern mit den Re\SourceLeaf{564}ductionselementen zusammenhängenden Umstand vorgenommen ist, hat man jedesmal die nöthigen Berichtigungen oder die Modificationen der Reductionselemente angebracht. Während der ersten Monate haben verschiedene hinlänglich geübte Beobachter sich mit mir in die Beobachtungen getheilt; vom 1. October 1834 an aber sind sie regelmässig durch Hrn.\ Doctor \textsc{Goldschmidt} angestellt, der nur in Behinderungsfällen durch andere geschickte Beobachter vertreten ist.

Die monatlichen Mittel aus diesen Bestimmungen bis Januar 1835 habe ich bereits in den Göttingischen gelehrten Anzeigen 1834 u.\ 1835 [S.\ 519 u.\ S.\ 528 d.\ B.] mitgetheilt: hier folgen nunmehr dieselben für drei vollständige Jahrgänge.

\Needspace{31\baselineskip}
\begin{center}\textit{Mittelwerth der westlichen magnetischen Declination zu Göttingen.}\end{center}
\begin{center}\small
\begin{tabular}{rl|rrr|rrr}
&&\multicolumn{3}{c|}{8 Uhr Vorm.}&\multicolumn{3}{c}{1 Uhr Nachm.}\\\hline
1834&März zweite Hälfte&$18^0$&$38'$&$16''0$&$18^0$&$46'$&$40''4$\\
&April&&36&6.9&&47&3.8\\
&Mai&&36&28.2&&47&15.4\\
&Junius&&37&40.7&&47&59.5\\
&Julius&&37&57.5&&48&19.0\\
&August&&38&48.1&&49&11.0\\
&September&&36&58.4&&46&32.3\\
&October&&37&18.4&&44&47.2\\
&November&&37&38.4&&43&4.3\\
&December&&37&54.8&&41&32.7\\[4pt]
1835&Januar&&37&51.5&&42&14.4\\
&Februar&&37&3.5&&42&29.4\\
&März&&34&47.5&&44&55.2\\
&April&&32&57.7&&46&31.6\\
&Mai&&32&13.4&&45&17.1\\
&Junius&&32&56.4&&44&41.3\\
&Julius&&34&8.0&&44&42.8\\
&August&&34&12.4&&46&56.8\\
&September&&33&21.2&&44&27.6\\
&October&&33&23.0&&43&5.3\\
&November&&36&15.3&&43&49.5\\
&December&&35&25.9&&40&19.1\\[4pt]
1836&Januar&&35&2.4&&40&34.6\\
&Februar&&33&26.7&&41&15.2\\
&März&&31&1.4&&43&16.4\\
&April&&26&32.9&&43&42.6\\
&Mai&&28&0.8&&44&37.2
\end{tabular}
\end{center}
\SourceLeaf{565}\Needspace{15\baselineskip}
\begin{center}\small
\begin{tabular}{rl|rrr|rrr}
&&\multicolumn{3}{c|}{8 Uhr Vorm.}&\multicolumn{3}{c}{1 Uhr Nachm.}\\\hline
1836&Junius&$18^0$&$27'$&$35''1$&$18^0$&$42'$&$52''4$\\
&Julius&&26&54.2&&42&26.0\\
&August&&25&42.4&&41&45.0\\
&September&&26&14.6&&40&59.6\\
&October&&27&34.0&&40&32.8\\
&November&&29&21.0&&36&54.3\\
&December&&29&13.7&&35&46.8\\[4pt]
1837&Januar&&27&35.3&&37&46.2\\
&Februar&&27&35.6&&36&28.3\\
&März&&25&44.2&&39&4.2
\end{tabular}
\end{center}

Es mögen nun einige Combinationen dieser Beobachtungen hier Platz finden.

Der Unterschied der Vormittags- und Nachmittags-Declination hat in den Mittelzahlen durchgängig einerlei Zeichen; die Abhängigkeit der Grösse dieses Unterschiedes von der Jahreszeit erkennt man in folgender Übersicht:
\begin{center}\small
\begin{tabular}{l||rr|rr|rr||rr}
&\multicolumn{2}{c|}{1834. 1835}&\multicolumn{2}{c|}{1835. 1836}&\multicolumn{2}{c||}{1836. 1837}&\multicolumn{2}{c}{Mittel}\\\hline\hline
April&$10'$&$56''9$&$13'$&$33''9$&$17'$&$9''7$&$13'$&$53''5$\\
Mai&10&47.2&13&3.7&16&36.4&13&29.1\\
Junius&10&18.8&11&44.9&15&17.3&12&27.0\\
Julius&10&21.5&10&34.8&15&31.8&12&9.4\\
August&10&22.9&12&44.4&16&2.6&13&3.3\\
September&9&33.9&11&6.4&14&45.0&11&48.4\\
October&7&28.8&9&42.3&12&58.8&10&3.3\\
November&5&25.9&7&34.2&7&33.3&6&51.1\\
December&3&37.9&4&53.2&6&33.1&5&1.4\\
Januar&4&22.9&5&32.2&10&10.9&6&42.0\\
Februar&5&25.9&7&48.5&8&52.7&7&22.4\\
März&10&7.7&12&15.0&13&20.0&11&54.2\\\hline\hline
Mittel&$8'$&$14''2$&$10'$&$2''8$&$12'$&$54''3$&$10'$&$23''8$
\end{tabular}
\end{center}
Man sieht, dass nicht blos in den Mittelwerthen, sondern auch in jedem einzelnen Jahre der Unterschied im December am kleinsten gewesen ist, und findet dies auch sehr natürlich, da die nach den Tageszeiten wechselnden Änderungen nothwendig einer Einwirkung der Sonne zugeschrieben werden müssen, wenn wir auch für jetzt noch nicht wissen, \textit{wie} diese Einwirkung geschieht. Dass dagegen die in den Sommermonaten ungleich grössern Unterschiede nicht um die
\SourceLeaf{566}Zeit des Solstitium am grössten, sondern im Junius und Julius kleiner waren, als im April, Mai und August, kann anfangs auffallend scheinen, zumal da die Übereinstimmung aller drei einzelnen Jahre in diesem Umstande eine Präsumtion gibt, dass dies nicht zufällig ist. Indessen darf dabei nicht übersehen werden, dass in den dem Solstitium nächsten Monaten die Zeit des Minimum der Declination schon auf eine frühere Stunde trifft, und daher die ganze Zunahme merklich grösser sein würde, als die Bewegung von 8 Uhr an gerechnet.

Es ist ferner auffallend, dass der Unterschied im zweiten Jahre in allen einzelnen Monaten grösser gewesen ist, als im ersten, und im dritten wieder grösser als im zweiten. Aber die Unterschiede sind viel zu gross, als dass man hierin etwas auf eine Säcularzunahme hinauslaufendes suchen dürfte, und es steht vielmehr zu erwarten, dass bei der Fortsetzung der Beobachtungen durch mehrere Jahre ein Hinundherschwanken nicht ausbleiben werde. Aber jedenfalls lernen wir daraus, dass auch bei dem Einwirken der Sonne auf den Erdmagnetismus ein Jahr vor dem andern ausgezeichnet sein kann, etwa ebenso, wie ein ganzer Sommer oder ein ganzer Winter von andern durch die Witterungsbeschaffenheit bedeutend verschieden ist. Eben deshalb aber wird man zu einer genauen Bestimmung der Mittelwerthe erst durch mehrjährige Beobachtungen gelangen können.

Dass ausnahmsweise an einzelnen Tagen der Unterschied der vormittägigen und nachmittägigen Declination das entgegengesetzte Zeichen haben kann, ist schon oben bemerkt. Die Seltenheit solcher Ausnahmen erhellt daraus, dass während der dreijährigen Beobachtungen nur vierzehn Fälle der Art vorgekommen sind, mithin durchschnittlich unter 79 Tagen einer. Ich setze sie hier her, nebst der Angabe, wie viel jedesmal die Declination 8 Uhr Morgens grösser gewesen ist, als 1 Uhr Nachmittags.
\begin{center}\small
\begin{tabular}{rlr|rr||rlr|rr}
1834&August&15&$6'$&$8''0$&1835&November&8&$3'$&$42''2$\\
&December&24&3&43.0&&December&8&18&35.6\\
&December&25&0&38.2&1836&Januar&20&0&46.3\\
&December&26&2&20.3&&Julius&20&5&8.8\\
1835&Januar&30&0&23.8&&November&9&11&9.5\\
&Februar&7&0&32.5&1837&Februar&13&4&1.0\\
&October&4&0&43.1&&März&14&1&22.6
\end{tabular}
\end{center}
Dass von diesen vierzehn Ausnahmen zwölf auf die Wintermonate und nur zwei auf die Sommermonate fallen, ist ganz in der Ordnung, da die geringe re\SourceLeaf{567}gelmässige Sonnenwirkung in den erstern leichter durch eine anomalische Bewegung überragt werden kann, als die viel grössere in den letztern.

Um zu versuchen, in wie fern sich aus den vorliegenden Beobachtungen die Säcularänderung schon erkennen lasse, sind die monatlichen Mittel des ersten Jahrs mit den entsprechenden des zweiten, und eben so die des zweiten mit denen des dritten verglichen. Unter den 48 auf diese Art hervorgehenden Vergleichungen (denn der unvollständige März 1834 ist von dieser wie von den übrigen Combinationen ausgeschlossen) geben 47 eine Abnahme, und nur eine eine Zunahme, welche deshalb in folgender Übersicht mit dem Minuszeichen bezeichnet ist.
\Needspace{17\baselineskip}
\begin{center}\small\setlength{\tabcolsep}{3pt}
\textit{Jährliche Abnahme der Declination.}\par\smallskip
\begin{tabular}{l||rr|rr||rr|rr||rr}
&\multicolumn{4}{c||}{Erstes Jahr}&\multicolumn{4}{c||}{Zweites Jahr}&\multicolumn{2}{c}{Mittel}\\
&\multicolumn{2}{c|}{8 Uhr Vorm.}&\multicolumn{2}{c||}{1 Uhr Nachm.}&\multicolumn{2}{c|}{8 Uhr Vorm.}&\multicolumn{2}{c||}{1 Uhr Nachm.}&&\\\hline\hline
April&$3'$&$9''2$&$0'$&$32''2$&$6'$&$24''8$&$2'$&$49''0$&$3'$&$13''8$\\
Mai&4&14.8&1&58.3&4&12.6&0&39.9&2&46.4\\
Junius&4&44.3&3&18.2&5&21.3&1&48.9&3&48.1\\
Julius&3&49.5&3&36.2&7&13.8&2&16.8&4&14.1\\
August&4&35.7&2&14.2&8&30.0&5&11.8&5&7.9\\
September&3&37.2&2&4.7&7&6.6&3&28.0&4&4.1\\
October&3&55.4&1&41.9&5&49.0&2&32.5&3&29.6\\
November&1&23.1&$-0$&45.2&6&54.3&6&55.2&3&36.8\\
December&2&28.9&1&13.6&6&12.2&4&32.3&3&36.7\\
Januar&2&49.1&1&39.8&7&27.1&2&48.4&3&41.1\\
Februar&3&36.8&1&14.2&5&51.1&4&46.9&3&52.2\\
März&3&46.1&1&38.8&5&17.2&4&12.2&3&43.6\\\hline\hline
Mittel&3&30.8&1&42.2&6&21.7&3&20.2&3&46.2
\end{tabular}
\end{center}
Dass die Vergleichung der vormittägigen Mittel hier meistens eine stärkere Abnahme gibt als die Vergleichung der nachmittägigen, ist nichts weiter als eine andere Einkleidung des schon oben bemerkten, dass die täglichen Änderungen im ersten Jahre geringer als im zweiten, und im zweiten geringer als im dritten gefunden waren. Es wird daher jener Unterschied nicht als ein reeller, sondern nur wie ein zufälliger zu betrachten, und bei längerer Fortsetzung der Beobachtungen auch ein Unterschied im entgegengesetzten Sinn zu erwarten sein. In so fern man also keinen hinreichenden Grund hat, dem einen Resultate vor dem andern einen Vorzug zu geben, bleibt nichts übrig, als sich an das Mittel aus
\SourceLeaf{568}beiden zu halten. Dieses Mittel ist beim ersten Jahre $2'36''5$, beim zweiten $4'55''9$, und man könnte versucht sein, dies als einen Beweis anzusehen, dass die Abnahme der Declination sich beschleunigt. Dies würde jedoch nichts weiter sein als ein schlechter Grund für eine an sich richtige Sache. Es ist nemlich bekannt, dass die während des vorigen Jahrhunderts in ganz Europa zunehmende Declination im gegenwärtigen ihr Maximum erreicht hat und seitdem wieder zurückgeht. Der Natur der Sache nach muss dieser Übergang eine anfangs unmerkliche und nach und nach stärker werdende Abnahme erzeugen. Allein obgleich in Ermangelung früherer Beobachtungen das Jahr, wo für Göttingen dieser Übergang Statt gefunden hat, sich nicht bestimmt angeben lässt, so muss man doch nach den von andern Orten bekannt gewordenen Beobachtungen dieses Jahr für beträchtlich weiter zurückliegend ansehen, als aus jenen beiden Zahlen folgen würde, wenn man sie als reine Wirkungen der langsamen Bewegung, die wir Säcularbewegung nennen, betrachten wollte. Und eben so ist nach allen sonstigen Erfahrungen eine so starke Änderung wie $2'19''4$ als regelmässige Zunahme für ein Jahr schlechterdings nicht zulässig. Wir halten daher auch diesen Unterschied grösstentheils für zufällig, so dass vor der Hand und bis weiter reichende Erfahrungen zu Gebote stehen werden, das Mittel $3'46''2$ als einjährige Abnahme der Declination für 1834---1837 gelten muss.

Da der Unterschied der Declinationen für die Vormittags- und die Nachmittagsstunde einer so offenbar mit der Jahreszeit wechselnden Ungleichheit unterworfen ist, so entsteht die Frage, ob nur die eine allein oder vorzugsweise, oder ob beide zugleich an einem von der Jahreszeit abhängenden Wechsel Theil nehmen, und welche Gesetze dabei zum Grunde liegen. Zur Ausmittelung dieser Gesetze wird zwar eine längere Reihe von Jahren noch nothwendiger sein, als für den blossen Unterschied der Declinationen: inzwischen wird man doch gern sehen, was die bisherigen Beobachtungen, so weit sie reichen, aussagen.

Es sind in dieser Absicht zuvörderst die Mittelwerthe aus je zwölf Monaten für die drei Beobachtungsjahre berechnet. Diese sind:
\begin{center}\small
\begin{tabular}{l|rrr|rrr}
&\multicolumn{3}{c|}{8 Uhr Vorm.}&\multicolumn{3}{c}{1 Uhr Nachm.}\\\hline
1834---1835&$18^0$&$37'$&$12''5$&$18^0$&$45'$&$27''0$\\
1835---1836&&33&42.0&&43&44.8\\
1836---1837&&27&20.3&&40&14.6
\end{tabular}
\end{center}

\SourceLeaf{569}Diese Mittelwerthe sind als gültig für den mittleren Tag jedes Rechnungsjahrs zu betrachten, also die ersten für den 1. October 1834 u. s. f.

Die Vergleichung der einzelnen Monate jedes Jahres mit dem zugehörigen Mittelwerthe gibt folgende Unterschiede:
\begin{center}\small
\textit{Declination 8 Uhr Vormittags.}\par\smallskip
\begin{tabular}{l|rr|rr|rr||rr}
&\multicolumn{2}{c|}{Erstes Jahr}&\multicolumn{2}{c|}{Zweites Jahr}&\multicolumn{2}{c||}{Drittes Jahr}&\multicolumn{2}{c}{Mittel}\\\hline
April&$-1'$&$5''9$&$-0'$&$44''3$&$-0'$&$47''4$&$-0'$&$52''5$\\
Mai&$-0$&44.6&$-1$&28.6&$+0$&40.5&$-0$&30.9\\
Junius&$+0$&27.9&$-0$&45.6&$+0$&14.8&$-0$&1.0\\
Julius&$+0$&44.7&$+0$&26.0&$-0$&26.1&$+0$&14.9\\
August&$+1$&35.3&$+0$&30.4&$-1$&37.9&$+0$&9.3\\
September&$-0$&14.4&$-0$&20.8&$-1$&5.7&$-0$&33.6\\
October&$+0$&5.6&$-0$&19.0&$+0$&13.7&$-0$&0.1\\
November&$+0$&25.6&$+2$&33.3&$+2$&0.7&$+1$&39.9\\
December&$+0$&42.0&$+1$&43.9&$+1$&53.4&$+1$&26.4\\
Januar&$+0$&38.7&$+1$&20.4&$+0$&15.0&$+0$&44.7\\
Februar&$-0$&9.3&$-0$&15.3&$+0$&15.3&$-0$&3.1\\
März&$-2$&25.3&$-2$&40.6&$-1$&36.1&$-2$&14.0
\end{tabular}
\end{center}
\Needspace{16\baselineskip}
\begin{center}\small
\textit{Declination 1 Uhr Nachmittags.}\par\smallskip
\begin{tabular}{l|rr|rr|rr||rr}
&\multicolumn{2}{c|}{Erstes Jahr}&\multicolumn{2}{c|}{Zweites Jahr}&\multicolumn{2}{c||}{Drittes Jahr}&\multicolumn{2}{c}{Mittel}\\\hline
April&$+1'$&$36''8$&$+2'$&$46''8$&$+3'$&$28''0$&$+2'$&$37''2$\\
Mai&$+1$&48.4&$+1$&32.3&$+4$&22.6&$+2$&34.4\\
Junius&$+2$&32.5&$+0$&56.5&$+2$&37.8&$+2$&2.3\\
Julius&$+2$&52.0&$+0$&58.0&$+2$&11.4&$+2$&0.5\\
August&$+3$&44.0&$+3$&12.0&$+1$&30.4&$+2$&48.8\\
September&$+1$&5.3&$+0$&42.8&$+0$&45.0&$+0$&51.0\\
October&$-0$&39.8&$-0$&39.5&$+0$&18.2&$-0$&20.4\\
November&$-2$&22.7&$+0$&4.7&$-3$&20.3&$-1$&52.8\\
December&$-3$&54.3&$-3$&25.7&$-4$&27.8&$-3$&55.9\\
Januar&$-3$&12.6&$-3$&10.2&$-2$&28.4&$-2$&57.1\\
Februar&$-2$&57.6&$-2$&29.6&$-3$&46.3&$-3$&4.5\\
März&$-0$&31.8&$-0$&28.4&$-1$&10.4&$-0$&43.5
\end{tabular}
\end{center}
Die Zahlen der letzten Columne sind als Mittel aus drei Jahren einigermaassen, wenn auch nur erst sehr unvollkommen, von dem Einflusse der unregelmässigen Anomalien befreit, allein offenbar noch mit der Säcularänderung behaftet. Um diese abzulösen, muss noch der Betrag derselben zwischen der
\SourceLeaf{570}Mitte jedes Monats und dem 1. October für die ersten sechs Monate mit negativem, für die letzten sechs mit positivem Zeichen angebracht werden. Unter Zugrundlegung des oben bestimmten zwölfmonatlichen Werths $3'46''2$ erhalten wir so folgende Resultate.
\begin{center}\small
\begin{tabular}{l|rr|rr||rr}
&\multicolumn{2}{c|}{8 Uhr Vorm.}&\multicolumn{2}{c||}{1 Uhr Nachm.}&\multicolumn{2}{c}{Mittel}\\\hline
April&$-2'$&$35''6$&$+0'$&$54''2$&$-0'$&$50''7$\\
Mai&$-1$&55.3&$+1$&10.0&$-0$&22.6\\
Junius&$-1$&6\phantom{.}6&$+0$&56.7&$-0$&4.9\\
Julius&$-0$&32.0&$+1$&13.6&$+0$&20.8\\
August&$-0$&18.8&$+2$&20.7&$+1$&0.9\\
September&$-0$&43.0&$+0$&41.6&$-0$&0.7\\
October&$+0$&9.3&$-0$&11.0&$-0$&0.8\\
November&$+2$&8.0&$-1$&24.7&$+0$&21.6\\
December&$+2$&13.3&$-3$&9.0&$-0$&27.8\\
Januar&$+1$&50.3&$-1$&51.5&$-0$&0.6\\
Februar&$+1$&21.3&$-1$&40.1&$-0$&9.4\\
März&$-0$&30.9&$+0$&59.6&$+0$&14.3
\end{tabular}
\end{center}
In diesen Resultaten zeigt sich schon so viele Regelmässigkeit, wie man von nur dreijährigen Beobachtungen erwarten konnte. Die erste Columne zeigt, wie viel die vormittägige Declination in den einzelnen Monaten von der mittlern vormittägigen Declination abweicht, und eben so gibt die zweite Columne den Unterschied der nachmittägigen Declination in jedem Monat von der mittlern nachmittägigen Declination, wobei man sich erinnern muss, dass die letztere selbst $10'23''8$ grösser ist, als die mittlere vormittägige.

Merkwürdig scheint nun, dass in allen zwölf Monaten die vormittägige und nachmittägige Declination \textit{auf entgegengesetzten Seiten} über ihre mittleren Werthe hinaus schwanken. In den fünf Wintermonaten von October bis Februar ist die vormittägige grösser als ihr mittlerer Werth, die nachmittägige kleiner, und beide Umstände tragen also \textit{zugleich} dazu bei, in dieser Jahreszeit die ganze Differenz unter ihren mittlern Werth zu bringen: in den übrigen sieben Monaten findet gerade das Entgegengesetzte Statt. Überdiess sind diese entgegengesetzten Schwankungen durchschnittlich nahe von gleicher Grösse, wovon die Folge ist, dass sie sich in ihrem Mittelwerth, welchen die letzte Columne darstellt, fast aufheben. Mit andern Worten ist dies auch so auszusprechen: das Mittel zwischen der magnetischen Declination Vormittags 8 Uhr und Nachmittags 1 Uhr
\SourceLeaf{571}enthält neben den unregelmässigen Anomalien und der Säcularabnahme keine erheblichen von der Jahreszeit abhängigen Schwankungen, wenigstens tritt gar kein Unterschied der Sommermonate gegen die Wintermonate mit Sicherheit hervor.

Der mittlere Werth selbst, aus sämmtlichen dreijährigen Beobachtungen abgeleitet, würde für den 1. October 1835
\[
=18^0\;37'\;56''9
\]
anzusetzen sein. Übrigens versteht sich von selbst, dass hier nur der Mittelwerth aus den bei unserm Beobachten gewählten Stunden gemeint ist, von welchem der Mittelwerth aus \textit{allen} Stunden des Tages wohl etwas verschieden sein könnte, wenn gleich wahrscheinlich nur wenig. Allein alle bisherigen Untersuchungen zeigen zur Genüge, dass ohne sehr langwierige Arbeiten darüber mit Sicherheit nichts wird festgesetzt werden können.

Bisher ist nur von den monatlichen Mittelzahlen die Rede gewesen. Der vollständige Abdruck der einzelnen Beobachtungen wurde für jetzt für überflüssig gehalten, da dieselben, so lange sie nur von Einem Orte vorliegen, nur in so fern ein Interesse haben könnten, als das unregelmässige Hinundherspringen sich daran erkennen lässt. Dieser Zweck lässt sich jedoch besser, als durch den blossen Anblick der Zahlen, vermittelst einer methodischen Combination derselben erreichen, wodurch die Grösse des Schwankens auf ein bestimmtes Maass zurückgeführt, und der allgemeine Charakter verschiedener Zeiträume, in Beziehung auf stärkeres oder geringeres Schwanken während derselben, genau vergleichbar wird. Ich verstehe hier Kürze halber unter dem Schwanken der magnetischen Declination die Differenz von der des vorhergehenden Tages zu derselben Stunde, und (nach der Analogie der sogenannten mittlern Beobachtungsfehler) unter mittlerm Schwanken während eines beliebigen Zeitraumes die Quadratwurzel aus dem Mittel der Quadrate der einzelnen Schwankungen. Man hat dabei zu bemerken, dass wenn mehrere gleiche oder als gleich betrachtete Zeiträume nachher zu einem einzigen vereinigt werden sollen, man zur Bestimmung des Generalmittels nicht das arithmetische Mittel aus den partiellen mittlern Schwankungen nehmen darf, sondern erst von letztern auf ihre Quadrate zurückkommen, aus diesen das arithmetische Mittel suchen muss, und sich an dessen Quadratwurzel zu halten hat. Die Resultate der auf diese Art über die dreijährigen
